\section*{Theorem 1 --- the even carrier code}
% =====================================================================
\addcontentsline{toc}{section}{Theorem 1}
Let $E$ be the rank-$\gcar=5$ carrier, split $E=E_3\oplus E_2$ (color $b=3$, weak $s=2$). The
one-generation spinor packet is the even exterior algebra, equivalently the even-weight binary code
\begin{equation}
  S^+=\Lambda^0E\oplus\Lambda^2E\oplus\Lambda^4E\ \cong\ C^+=\{x\in\mathbb F_2^5:|x|\ \text{even}\},\qquad
  \dim S^+=|C^+|=2^{\gcar-1}=16 .
\end{equation}
The two carrier roots $y_+=\tfrac12$, $y_-=-\tfrac13$ generate all hypercharges
($Y(Q_L)=y_++y_-=\tfrac16$, $Y(u^c)=2y_-$, $Y(e^c)=2y_+$, $Y(\nu^c)=0$) and the transport cusps
$\{1,\tfrac23,\tfrac13\}=\{2y_+,-2y_-,2(y_++y_-)\}$, with $\gamma_{\mathrm{car}}=\operatorname{Tr}_EY^2=3y_-^2+2y_+^2=\tfrac56$. \mE
\smallskip\\
\noindent\textbf{Gauge-rank factorisation (\veri{v79\_review\_identities.py}).} The carrier polynomial factors
into the two gauge ranks directly:
\[
  \boxed{\ 6Y^2-Y-1=(2Y-1)(3Y+1)\ },\qquad Y=\tfrac12\ (s{=}2,\ SU(2)),\quad Y=-\tfrac13\ (b{=}3,\ SU(3)),
\]
i.e.\ $(sY-1)(bY+1)$ --- the polynomial splits into one linear representative of each Standard-Model
non-abelian factor; the $5$-slot carrier \emph{forces} the split into $SU(3)\times SU(2)$.

\begin{keybox}[Hypercharge Lucas--Binet resonances \mE\ (mode reading \mC, \veri{v79\_review\_identities.py})]
The integer carrier charges $(q_+,q_-)=(3,-2)$ (the primitive trace-free pair, $X^2{-}X{-}6=(X{-}3)(X{+}2)$,
Lean \texttt{unique\_carrier\_pair}) generate the Lucas $U$-sequence
\[
  D_n=\frac{3^n-(-2)^n}{5}:\qquad D_1{,}\dots{,}D_6=1,\ 1,\ 7,\ 13,\ 55,\ 133,
\]
which is, \emph{in order}, $(N_\Phi,\ \cdot\ ,\ \text{scalaron exponent }7,\ \Delta_Q{=}13,\ \text{quark numerator }55,\ \dim E_7{=}133)$.
In particular the quark ratio reads $c_u/c_d=D_5/(\Nfam^2 D_4)=\tfrac{55}{117}$ --- a single recurrence of the
hypercharge generator, not a free integer. \mE\ for the arithmetic; the per-mode physical identification is a
\mC\ cross-domain reading (only $n\le6$, the carrier-slot range, is meaningful).
\end{keybox}

% =====================================================================
\section*{Theorem 2 --- Exterior--Action Duality}
% =====================================================================
\addcontentsline{toc}{section}{Theorem 2}
On $C^+$ the flip $F=x\oplus y$ is even, so $|F|\in\{0,2,4\}$. The three relations
$R_0,R_2,R_4$ form a symmetric association scheme (the even Hamming scheme on five bits) with
\begin{equation}
  \boxed{\ \begin{aligned}
    &\text{valencies}\ (1,10,5)=(\dim\Lambda^0E,\dim\Lambda^2E,\dim\Lambda^4E),\\
    &\text{spectral multiplicities}\ (1,5,10)=\AEW{:}\AH{:}\AL
  \end{aligned}\ }\qquad\mE .
\end{equation}
The eigenvalues (computed directly) are
$A_2:\{10^{(1)},2^{(5)},(-2)^{(10)}\}$ and $A_4:\{5^{(1)},(-3)^{(5)},1^{(10)}\}$, with
$A_2+A_4=K_{16}-\mathbf 1$. The two graphs are named objects:
\begin{equation}
\begin{aligned}
  &R_2(C^+)=\text{folded/complement, } \mathrm{SRG}(16,10,6,6)\ (\text{the Clebsch complement});\\
  &\boxed{\ R_4(C^+)=\text{Clebsch graph, } \mathrm{SRG}(16,5,0,2)\ }\qquad\mE .
\end{aligned}
\end{equation}
\begin{remark}
The Standard-Model family is the \emph{valency} side $(1,10,5)$ of one scheme; the cosmological action
ladder is its \emph{spectral} side $(1,5,10)$. The action ladder is therefore not merely ``also
binomial'' --- it is the spectral dual of the even exterior family. \mC
\end{remark}
\noindent\textbf{This duality is the MacWilliams transform (\veri{v79\_review\_identities.py}).} $C^+$ is the
even-weight code $[5,4,2]$; its dual $(C^+)^\perp$ is the repetition code $[5,1,5]=\{00000,11111\}$ (the
uniform/vacuum word). The MacWilliams identity sends the weight enumerator of the dual to that of $C^+$, namely
$(1,10,5)$ --- so the valency$\leftrightarrow$spectral duality \emph{is} the coding-theoretic UV/IR duality, and
the carrier grading $\Lambda^{\mathrm{even}}(5)=1{+}10{+}5$ is the MacWilliams image of the repetition vacuum.
\mE\ (code identity) / \mC\ (holographic reading).

\noindent\textbf{The full Bose--Mesner algebra} (machine-checkable) fixes the duality beyond
eigenvalues:
\begin{equation}
  A_2^2=10\,\mathbf 1+6A_2+6A_4,\quad A_2A_4=3A_2+4A_4,\quad A_4^2=5\,\mathbf 1+2A_2 ,
\end{equation}
with eigenmatrix and multiplicity orthogonality
\begin{equation}
  P=\begin{pmatrix}1&10&5\\1&2&-3\\1&-2&1\end{pmatrix},\quad \det P=-64,\qquad
  \boxed{\ P^{\!\top}\operatorname{diag}(1,5,10)\,P=\operatorname{diag}(16,160,80)\ }\qquad\mE ,
\end{equation}
so the flip split $240=160+80$ ($|F|=2,4$) appears directly in the scheme's eigenmatrix orthogonality.

\noindent\textbf{Carrier $=D_5$ (verified).} The $16$ vertices of $C^+$ are exactly the weights of the
$D_5=\mathfrak{so}_{10}$ half-spinor $\tfrac12(\pm1,\dots,\pm1)$ with an even number of minus signs.
The flip-2 and flip-4 graphs are translation- and permutation-invariant, so the \emph{affine} group
$C^+\rtimes S_5$ (translations $\rtimes$ coordinate permutations) acts by graph/scheme automorphisms; the
group closure gives the automorphism group of the even Hamming scheme (equivalently the Clebsch graph),
\begin{equation}
  \boxed{\ \operatorname{Aut}_{\mathrm{aff}}(C^+)=C^+\rtimes S_5,\qquad |\operatorname{Aut}_{\mathrm{aff}}(C^+)|=2^{\gcar-1}\,\gcar!=16\cdot120=1920=|W(D_5)|\ }\qquad\mE ,
\end{equation}
and $1920=r\,|R(E_8)|=8\cdot240$. (The \emph{linear} code automorphism group alone is only the
coordinate group $S_5$ of order $120$; the factor $2^{\gcar-1}=16$ is the translation/scheme part. It is
the affine/scheme group that equals $W(D_5)$.) Edges of the Clebsch graph: $|E(R_4)|=\tfrac{16\cdot5}{2}=40=|R(D_5)|$
$=\sum_{f,j}L_{f,j}$ (the entire flavor transport budget). So the five-slot even code is the half-spin
weight structure of $\mathfrak{so}_{10}$, the classical home of one SM family with $\nu^c$.

% =====================================================================
\section*{Theorem 3 --- the Coxeter closure equation selects $\gcar=5$}
% =====================================================================
\addcontentsline{toc}{section}{Theorem 3}
Requiring the even family to equal the action truncation through pair degree,
$\dim S^+=2^{g-1}=\binom g0+\binom g1+\binom g2$, i.e.
\begin{equation}
  \boxed{\ 2^{g}-2=g(g+1)\ \Longleftrightarrow\ g=5\ }\qquad\mE,
\end{equation}
the unique positive-integer solution. At $g=5$ the two sides equal $30$, so
\emph{carrier exhaustion} $2^{g}-2$ and the \emph{Coxeter period} $g(g+1)$ coincide only at five
slots: $2^{\gcar}-2=30=h$ and $\gcar(\gcar+1)=30$.

% =====================================================================
\section*{Theorem 4 --- the $\mathbb Z_{30}$ cyclotomic Coxeter compiler}
% =====================================================================
\addcontentsline{toc}{section}{Theorem 4}
From the three atoms $(2,3,\gcar)$, with $h=2\cdot3\cdot\gcar=30$ and $r=\varphi(h)=8$:
\begin{equation}
  \boxed{\ |R(E_8)|=rh=240,\qquad \dim E_8=r(h+1)=8\cdot31=248,\qquad h+1=2^{\gcar}-1=31\ }\qquad\mE .
\end{equation}
The $E_8$ exponents are the primitive residues modulo $30$,
\begin{equation}
  \operatorname{Exp}(E_8)=(\mathbb Z/30\mathbb Z)^\times=\{1,7,11,13,17,19,23,29\},\qquad
  \chi_C(t)=\Phi_{30}(t)=t^8+t^7-t^5-t^4-t^3+t+1 ,
\end{equation}
and $|(\mathbb Z/30)^\times|=\varphi(30)=8=r$. The local SM algebra and the root atlas are the two
arithmetic projections of $(r,h)=(8,30)$:
\begin{equation}
  \Oadm=6r=48,\quad D_{\mathrm{start}}=2h=60,\quad \gcd(6r,2h)=\dimg=12,\quad \operatorname{lcm}(6r,2h)=rh=240 .
\end{equation}

% =====================================================================
\section*{Theorem 5 --- the root transition atlas and the positive-edge moment}
% =====================================================================
\addcontentsline{toc}{section}{Theorem 5}
The directed even-flips of $C^+$ are the complete graph $K_{16}$:
\begin{equation}
  \boxed{\ \begin{gathered}
  |R^+(E_8)|=\binom{\dim S^+}{2}=\binom{16}{2}=120=\operatorname{Tr}_{S^+}X^2=5!=\!\!\sum_{m\in(\mathbb Z/30)^\times}\!\!\!m,\\[2pt]
  |R(E_8)|=2\binom{16}{2}=240\ \end{gathered}}\quad\mE .
\end{equation}
($X=6Y$.) So the positive $E_8$ roots are the unordered transitions of one family; the second
hypercharge moment, the positive-root count and the sum of the $E_8$ exponents coincide at $120$.
The CP pairing $m\leftrightarrow 30-m$ of the exponents gives four pairs, each of sum $h=30$:
\begin{equation}
  \boxed{\ \operatorname{Tr}_{S^+}X^2=(N_{\mathrm{fam}}+N_\Phi)\,h=4\cdot30=120=\dim\mathfrak{so}_{16}=\underbrace{45}_{D_5}+\underbrace{15}_{A_3}+\underbrace{60}_{(10,6)}\ }\qquad\mE ,
\end{equation}
so the hypercharge moment is simultaneously four CP-paired Coxeter periods \emph{and} the adjoint
half $\mathfrak{so}_{16}=(45,1)\oplus(1,15)\oplus(10,6)$ of the $E_8$ closure.
\par\smallskip\noindent\textbf{How the root \emph{origin} is now obtained --- see Theorem 7.}
A naive flip-size-preserving map $R_{\mathrm{flip}}\to\mathbb R^8$ cannot be $E_8$: the flip split
$240=160+80$ ($|F|=2,4$) differs from the $D_8$ split $112+128$. The correct route is \emph{not} a
direct $240$-metrization but the standard branching $E_8\supset D_5\times A_3$ (Theorem 7), which
builds the root spaces from known $D_5$ and $A_3$ data. Until the $A_3$ lattice is derived from the
four punctures (Theorem 7, \mO), $E_8$ remains a Coxeter atlas with a concrete, falsifiable lift.
\par\smallskip

% =====================================================================
\section*{Theorem 6 --- downstream readouts from $(h,r,\cthree)$}
% =====================================================================
\addcontentsline{toc}{section}{Theorem 6}
All discrete budgets and the seam seed are functions of $(h,r,\cthree)=(30,8,1/(8\pi))$:
\begin{align}
  \dim S^+&=2r=16, & h&=2(\dim S^+-1)=30, \\
  10b_1&=5r+1=41=1+h\tfrac{\dim S^+}{\dimg}, & I_1^{\mathrm{EW}}&=\frac{5r+1}{6r}=\frac{41}{48}, \\
  \useed&=\frac{r}{6}\cthree+6r\,\cthree^4=\frac43\cthree+48\cthree^4, &
  \dtwo&=4!\,\tbinom{16}{2}\,\cthree^8=2880\,\cthree^8, \\
  \xi_{\mathrm{tree}}&=\frac{6}{r}=\frac{\dimg}{\dim S^+}=\frac34, &
  \frac{M_{\mathrm{scal}}^2}{\Mbar^2}&=\cthree^{\,\Oadm-10b_1}=\cthree^{\,r-1}=\cthree^7 .
\end{align}
\textbf{Two-graph product table.} With the carrier-pair graph $K_5$ ($|E(K_5)|=\binom52=10=\AL$) and
the family-corner graph $K_4$ on $\mu_4$ ($|E(K_4)|=\binom42=6$, $|R(A_3)|=2|E(K_4)|=12$), every large
budget is a product of the two graphs:
\begin{equation}
  \boxed{\ 40=10\cdot4,\quad 60=10\cdot6,\quad 120=10\cdot12,\quad 240=10\cdot24\ }\qquad\mE ,
\end{equation}
i.e.\ ($\sum L$, $D_{\mathrm{start}}$, $|R^+(E_8)|$, $|R(E_8)|$) $=10\cdot(|\mu_4|,|E(K_4)|,|R(A_3)|,|S_4|)$.
Two occupancy identities follow,
\begin{equation}
\begin{aligned}
  &\boxed{\ 10b_1=1+|\mu_4|\,|E(K_5)|=1+4\cdot10=41\ },\\
  &\boxed{\ \Oadm=\Nfam\dim S^+=(\Nfam{+}N_\Phi)|R(A_3)|=3\cdot16=4\cdot12=48\ }\qquad\mE ,
\end{aligned}
\end{equation}
so the Higgs singleton is the fourth ($\mu_4$) direction that closes $\Nfam=3$ to $\mathbf 4$ of $A_3$.
\textbf{Occupancy deficit ladder.} The admissible occupancy splits cleanly,
\begin{equation}
  \boxed{\ \begin{aligned}
    &\Oadm=\underbrace{40}_{\sum L\,(\text{flavor})}+\underbrace{1}_{\text{Higgs}}+\underbrace{7}_{\text{scalaron}}=48,\qquad
    \Oadm-\!\sum L=8=h(D_5)=r(E_8),\\
    &\Oadm-10b_1=7=h(D_5)-1
  \end{aligned}\ }\qquad\mE ,
\end{equation}
so the scalaron exponent $r-1=7$ is exactly the post-flavor, post-Higgs occupancy remainder,
$M_{\mathrm{scal}}^2/\Mbar^2=\cthree^{\,\Oadm-10b_1}=\cthree^7$. With $h(D_5)=8$ and $|R^+(A_3)|=6$
the seed itself is a $D_5/A_3$ ratio--product, $\useed=\tfrac{h(D_5)}{|R^+(A_3)|}\cthree+h(D_5)|R^+(A_3)|\cthree^4$,
and $\xi=\tfrac{|R^+(A_3)|}{h(D_5)}\,(1+|R^+(A_3)|^2\cthree^3)^{-1}=\tfrac34(1+36\cthree^3)^{-1}$.
The seed readouts and the two decuples are
\begin{equation}
  \lC^2=\useed(1-\useed),\quad \betarad=2\cthree\useed,\quad \sin^2\theta_{13}=e^{-5/6}\useed,\quad
  \boxed{\ \omega_b=\frac{\lC}{\AL}=\frac{\lC}{10}=0.02244\ }\ (\text{Planck }0.02237),
\end{equation}
\begin{equation}
  \boxed{\ \begin{aligned}
    &\det M_\ell\propto\useed^{\,h/\Nfam}=\useed^{10},\\
    &\frac{\rho_\Lambda}{\Mbar^4}=\frac{3}{4\pi^2}e^{2\dph}\Bigl[\tfrac{v_{\mathrm{geo}}}{5\betarad^2\Mbar}\Bigr]^{h/\Nfam}
  \end{aligned}\ },
\end{equation}
with lepton powers $(5,3,2)=(\gcar,b,s)$: $\sum=10=\AL$ (the $\Lambda$ decuple), $\prod=30=h$ (the
Coxeter period). The first-generation transport sums read the non-primitive $\mathbb Z_{30}$ load,
\begin{equation}
  (G_1,G_2,G_3)=(22,13,5)=(h-r,\ \dimg+1,\ \gcar),\qquad h-r=30-\varphi(30)=22 ,
\end{equation}
\textbf{Flavor transport is a $D_5\times A_3$ object.} The transport packets
$L_u{=}(7,3,0),L_d{=}(7,5,2),L_e{=}(8,5,3)$ split into a $D_5$ budget and an $A_3$ geometry:
\begin{equation}
  \boxed{\ \textstyle\sum_{f,j}L_{f,j}=40=|R(D_5)|\ (\text{Clebsch edges}),\qquad
  L_{f,j}\bmod 6\ \text{lives on}\ C_6=Z_6=|R^+(A_3)|\ }\qquad\mE ,
\end{equation}
the residues being the $D_6$-orbit of the perfect set $\{0,1,3\}\subset C_6$ (cyclic distances
$\{1,2,3\}$): $u\!\equiv\!\{0,1,3\}$, $d\!\equiv\!2{-}\{0,1,3\}$, $e\!\equiv\!2{+}\{0,1,3\}$, with the
universal first-generation $+6$ winding. So the carrier $D_5$ sets the total budget and the family
$A_3$ (its $6$-element positive-root hexagon) sets the per-entry distances.
The scalaron mass is a pure seam power \mE; the amplitude is a \emph{conditional} readout, kept
\mC/\mO\ until the scalaron variable is formally identified with the Starobinsky $M$ parameter. That
identification is now \emph{heat-kernel grounded} (G2): the spectral-action Seeley--DeWitt coefficient
$a_4|_{R^2}\!\sim\!R^2/72$ delivers exactly the Starobinsky $R^2$ term, with $M^2/\Mbar^2=6(4\pi)^2/f_0$ and
the TFPT closure $\cthree^7$ fixing the cutoff moment $f_0$ (\veri{v36\_spectral\_action\_g2.py}):
\begin{equation}
  \underbrace{\frac{M_{\mathrm{scal}}^2}{\Mbar^2}=\cthree^{\,r-1}=\cthree^7}_{\mE},\qquad
  \underbrace{A_s=\frac{N_\star^2}{24\pi^2}\frac{M_{\mathrm{scal}}^2}{\Mbar^2}=\frac{N_\star^2\cthree^7}{24\pi^2}=2.17\times10^{-9}}_{\mC/\mO\ (\text{if }M_{\mathrm{scal}}=M_{\mathrm{Star}})} .
\end{equation}

% =====================================================================
